% GENERATED FILE. Edit canonical lessons, then run npm run generate:books.
% canonical-source-hash: fnv1a64:fb3ba477ecd73e57
% canonical-lessons: TE-C266-mallepuvvu, TE-C266-vepa, TE-C266-marri, TE-C266-tulasi, TE-C266-mogga, TE-R266-first-pass-words-from-chapters-258-261, TE-R266-second-pass-words-from-chapters-262-266

\chapter{Jasmine flower, Neem, Banyan tree, Holy basil, Bud}
\label{ch:te-jasmine-flower-neem-banyan-tree-holy-basil-bud}
\hlchaptermodality{266}

\begin{chapteropening}
\textbf{By the end of this chapter:} \emph{I can say a jasmine flower, neem, a banyan tree, holy basil and a bud.}

Complete the last lesson of chapter 266: I can say a jasmine flower, neem, a banyan tree, holy basil and a bud.
\end{chapteropening}

\section[\texorpdfstring{\te{మల్లెపువ్వు} (mallepuvvu)}{మల్లెపువ్వు (mallepuvvu)}]{\te{మల్లెపువ్వు} (mallepuvvu) — a jasmine flower}

\label{lesson:TE-C266-mallepuvvu}



\begin{quote}
Before the new one: say the Telugu for a thief, then the Telugu for an animal.
\end{quote}

\subsection*{You'll want to know: \te{మల్లెపువ్వు}}
\textbf{\te{మల్లెపువ్వు}} — \emph{mallepuvvu} — \textquotedblleft{}a jasmine flower\textquotedblright{}.

\begin{grammarlens}[title={What you've built}]
One more word for this chapter.
\end{grammarlens}

\subsection*{Your turn}
\begin{itemize}
\raggedright
  \item \emph{Say it:} \emph{mallepuvvu}
  \item \emph{Say it:} \emph{mallepuvvu}, once more
  \item \emph{Say it:} say \emph{mallepuvvu}
  \item \emph{Recall:} say the Telugu for a thief, then the Telugu for an animal, then say \emph{mallepuvvu} again
\end{itemize}

\begin{tcolorbox}[breakable,colback=teal!4,colframe=teal!35!black,title={Before you move on}]
What does \te{మల్లెపువ్వు} mean? (\textquotedblleft{}A jasmine flower\textquotedblright{}.) Say it once more.
\end{tcolorbox}
\section[\texorpdfstring{\te{వేప} (vēpa)}{వేప (vēpa)}]{\te{వేప} (vēpa) — neem, the neem tree}

\label{lesson:TE-C266-vepa}



\begin{quote}
Before the new one: say the Telugu for an animal, then the Telugu for a jasmine flower.
\end{quote}

\subsection*{You'll want to know: \te{వేప}}
\textbf{\te{వేప}} — \emph{vēpa} — \textquotedblleft{}neem, the neem tree\textquotedblright{}.

\begin{grammarlens}[title={What you've built}]
One more word for this chapter.
\end{grammarlens}

\subsection*{Your turn}
\begin{itemize}
\raggedright
  \item \emph{Say it:} \emph{vēpa}
  \item \emph{Say it:} \emph{vēpa}, once more
  \item \emph{Say it:} say \emph{vēpa}
  \item \emph{Recall:} say the Telugu for an animal, then the Telugu for a jasmine flower, then say \emph{vēpa} again
\end{itemize}

\begin{tcolorbox}[breakable,colback=teal!4,colframe=teal!35!black,title={Before you move on}]
What does \te{వేప} mean? (\textquotedblleft{}Neem, the neem tree\textquotedblright{}.) Say it once more.
\end{tcolorbox}
\section[\texorpdfstring{\te{మర్రి} (marri)}{మర్రి (marri)}]{\te{మర్రి} (marri) — a banyan tree}

\label{lesson:TE-C266-marri}



\begin{quote}
Before the new one: say the Telugu for a jasmine flower, then the Telugu for neem.
\end{quote}

\subsection*{You'll want to know: \te{మర్రి}}
\textbf{\te{మర్రి}} — \emph{marri} — \textquotedblleft{}a banyan tree\textquotedblright{}.

\begin{grammarlens}[title={What you've built}]
One more word for this chapter.
\end{grammarlens}

\subsection*{Your turn}
\begin{itemize}
\raggedright
  \item \emph{Say it:} \emph{marri}
  \item \emph{Say it:} \emph{marri}, once more
  \item \emph{Say it:} say \emph{marri}
  \item \emph{Recall:} say the Telugu for a jasmine flower, then the Telugu for neem, then say \emph{marri} again
\end{itemize}

\begin{tcolorbox}[breakable,colback=teal!4,colframe=teal!35!black,title={Before you move on}]
What does \te{మర్రి} mean? (\textquotedblleft{}A banyan tree\textquotedblright{}.) Say it once more.
\end{tcolorbox}
\section[\texorpdfstring{\te{తులసి} (tulasi)}{తులసి (tulasi)}]{\te{తులసి} (tulasi) — holy basil, tulsi}

\label{lesson:TE-C266-tulasi}



\begin{quote}
Before the new one: say the Telugu for neem, then the Telugu for a banyan tree.
\end{quote}

\subsection*{You'll want to know: \te{తులసి}}
\textbf{\te{తులసి}} — \emph{tulasi} — \textquotedblleft{}holy basil, tulsi\textquotedblright{}.

\begin{grammarlens}[title={What you've built}]
One more word for this chapter.
\end{grammarlens}

\subsection*{Your turn}
\begin{itemize}
\raggedright
  \item \emph{Say it:} \emph{tulasi}
  \item \emph{Say it:} \emph{tulasi}, once more
  \item \emph{Say it:} say \emph{tulasi}
  \item \emph{Recall:} say the Telugu for neem, then the Telugu for a banyan tree, then say \emph{tulasi} again
\end{itemize}

\begin{tcolorbox}[breakable,colback=teal!4,colframe=teal!35!black,title={Before you move on}]
What does \te{తులసి} mean? (\textquotedblleft{}Holy basil, tulsi\textquotedblright{}.) Say it once more.
\end{tcolorbox}
\section[\texorpdfstring{\te{మొగ్గ} (mogga)}{మొగ్గ (mogga)}]{\te{మొగ్గ} (mogga) — a bud}

\label{lesson:TE-C266-mogga}



\begin{quote}
Before the new one: say the Telugu for a banyan tree, then the Telugu for holy basil.
\end{quote}

\subsection*{You'll want to know: \te{మొగ్గ}}
\textbf{\te{మొగ్గ}} — \emph{mogga} — \textquotedblleft{}a bud\textquotedblright{}.

\begin{grammarlens}[title={What you've built}]
One more word for this chapter.
\end{grammarlens}

\subsection*{Your turn}
\begin{itemize}
\raggedright
  \item \emph{Say it:} \emph{mogga}
  \item \emph{Say it:} \emph{mogga}, once more
  \item \emph{Say it:} say \emph{mogga}
  \item \emph{Recall:} say the Telugu for a banyan tree, then the Telugu for holy basil, then say \emph{mogga} again
\end{itemize}

\begin{tcolorbox}[breakable,colback=teal!4,colframe=teal!35!black,title={Before you move on}]
What does \te{మొగ్గ} mean? (\textquotedblleft{}A bud\textquotedblright{}.) Say it once more.
\end{tcolorbox}
\section[(dialogue)]{Words from chapters 258-261 — a first pass}

\label{lesson:TE-R266-first-pass-words-from-chapters-258-261}



\begin{quote}
Say the words for holy basil and a bud.
\end{quote}

\subsection*{Your turn}
Say these aloud:

\begin{itemize}
\raggedright
  \item a method, a rule, an effect, development, a situation
  \item crowded, important, own, busy, compulsory
  \item smooth, warm, rare, strange, kind
  \item natural, thick, greasy, sticky, rough
\end{itemize}

\begin{tcolorbox}[breakable,colback=teal!4,colframe=teal!35!black,title={Before you move on}]
A method? (\textbf{\te{పద్ధతి}}, \emph{paddhati}.) A rule? (\textbf{\te{నియమం}}, \emph{niyamaṁ}.) An effect? (\textbf{\te{ప్రభావం}}, \emph{prabhāvaṁ}.) Development? (\textbf{\te{అభివృద్ధి}}, \emph{abhivṛddhi}.) A situation? (\textbf{\te{పరిస్థితి}}, \emph{paristhiti}.) Crowded? (\textbf{\te{రద్దీ}}, \emph{raddī}.) Important? (\textbf{\te{ముఖ్యమైన}}, \emph{mukhyamaina}.) Own? (\textbf{\te{సొంత}}, \emph{sonta}.) Busy? (\textbf{\te{బిజీ}}, \emph{bijī}.) Compulsory? (\textbf{\te{తప్పనిసరి}}, \emph{tappanisari}.) Smooth? (\textbf{\te{నునుపైన}}, \emph{nunupaina}.) Warm? (\textbf{\te{వెచ్చని}}, \emph{veccani}.) Rare? (\textbf{\te{అరుదైన}}, \emph{arudaina}.) Strange? (\textbf{\te{వింత}}, \emph{vinta}.) Kind? (\textbf{\te{దయగల}}, \emph{dayagala}.) Natural? (\textbf{\te{సహజమైన}}, \emph{sahajamaina}.) Thick? (\textbf{\te{చిక్కని}}, \emph{cikkani}.) Greasy? (\textbf{\te{జిడ్డు}}, \emph{jiḍḍu}.) Sticky? (\textbf{\te{జిగట}}, \emph{jigaṭa}.) Rough? (\textbf{\te{గరుకు}}, \emph{garuku}.)
\end{tcolorbox}
\section[(dialogue)]{Words from chapters 262-266 — a second pass}

\label{lesson:TE-R266-second-pass-words-from-chapters-262-266}



\begin{quote}
Say the words for bread and a chapati.
\end{quote}

\subsection*{Your turn}
Say these aloud:

\begin{itemize}
\raggedright
  \item bread, a chapati, upma, breakfast, green gram
  \item peanuts, taste, trousers, shorts, a vest
  \item a cheque, a visa, luggage, a newspaper, a number
  \item a bride, people, God, a thief, an animal
  \item a jasmine flower, neem, a banyan tree, holy basil, a bud
\end{itemize}

\begin{tcolorbox}[breakable,colback=teal!4,colframe=teal!35!black,title={Before you move on}]
Bread? (\textbf{\te{రొట్టె}}, \emph{roṭṭe}.) A chapati? (\textbf{\te{చపాతీ}}, \emph{capātī}.) Upma? (\textbf{\te{ఉప్మా}}, \emph{upmā}.) Breakfast? (\textbf{\te{టిఫిను}}, \emph{ṭiphinu}.) Green gram? (\textbf{\te{పెసలు}}, \emph{pesalu}.) Peanuts? (\textbf{\te{వేరుశనగ}}, \emph{vēruśanaga}.) Taste? (\textbf{\te{రుచి}}, \emph{ruci}.) Trousers? (\textbf{\te{ప్యాంటు}}, \emph{pyāṇṭu}.) Shorts? (\textbf{\te{నిక్కరు}}, \emph{nikkaru}.) A vest? (\textbf{\te{బనీను}}, \emph{banīnu}.) A cheque? (\textbf{\te{చెక్కు}}, \emph{cekku}.) A visa? (\textbf{\te{వీసా}}, \emph{vīsā}.) Luggage? (\textbf{\te{సామాను}}, \emph{sāmānu}.) A newspaper? (\textbf{\te{వార్తాపత్రిక}}, \emph{vārtāpatrika}.) A number? (\textbf{\te{నంబరు}}, \emph{nambaru}.) A bride? (\textbf{\te{పెళ్ళికూతురు}}, \emph{peḷḷikūturu}.) People? (\textbf{\te{జనం}}, \emph{janaṁ}.) God? (\textbf{\te{దేవుడు}}, \emph{dēvuḍu}.) A thief? (\textbf{\te{దొంగ}}, \emph{doṅga}.) An animal? (\textbf{\te{జంతువు}}, \emph{jantuvu}.) A jasmine flower? (\textbf{\te{మల్లెపువ్వు}}, \emph{mallepuvvu}.) Neem? (\textbf{\te{వేప}}, \emph{vēpa}.) A banyan tree? (\textbf{\te{మర్రి}}, \emph{marri}.) Holy basil? (\textbf{\te{తులసి}}, \emph{tulasi}.) A bud? (\textbf{\te{మొగ్గ}}, \emph{mogga}.)
\end{tcolorbox}
